\section{Gethsemane and the Open Account}

The unpayable debt of creation, the worth of obedience when feeling recoils, the sacrifice that heals rather than bargains, the freedom of submission, and the refusal to call suffering good in itself all converge in one night outside Jerusalem. In the Garden of Gethsemane, the question of the creature's due return receives its decisive human answer from the only Person who could offer it completely.

\subsection{The Created Human Will of the Incarnate Son}

St.\ Matthew and St.\ Luke record the scene without ornament. Entering Gethsemane with Peter, James, and John, Jesus ``began to grow sorrowful and to be sad,'' telling them, ``My soul is sorrowful even unto death'' (Matt 26:37--38). Withdrawing a stone's cast, he falls upon his face and kneels in agony until his sweat becomes as drops of blood:

\begin{quote}
My Father, if it be possible, let this chalice pass from me. Nevertheless, not as I will but as thou wilt. (Matt 26:39; cf.\ Luke 22:42: ``Father, if thou wilt, remove this chalice from me: but yet not my will, but thine be done.'')
\end{quote}

To read Gethsemane rightly requires strict Christological precision. The Person praying in the garden is the eternal, uncreated Son of God, consubstantial with the Father and sharing one undivided divine essence and one divine will with the Father and the Holy Spirit. Nothing in the divine Person is created. Yet in the Incarnation, that eternal Son truly assumed a complete, integral \emph{created} human nature---human body, human intellect, human affections, and a genuine created human will.

In the fourth century, St.\ Gregory of Nazianzus faced Arian opponents who cited Gethsemane and John 6:38 to claim that the Son had an inferior divine will subordinate to the Father's. In his Fourth Theological Oration (\emph{Oration} 30.12), Gregory fiercely guards the unity of the divine will---``as We have one Godhead, so We have one Will''---while recognizing that the shrinking expressed in Gethsemane belongs to our assumed human nature, ``inasmuch as the human will does not completely follow the Divine, but for the most part struggles against and resists it,'' whereas in the Saviour his human will ``cannot be opposed to God, seeing it is altogether taken into God.''\endnote{St.\ Gregory of Nazianzus, \emph{Oration} 30 (Fourth Theological Oration), \S 12 (Browne--Swallow translation, NPNF II.7).}

Three centuries later, against the Monothelites who denied a distinct human will in Christ altogether, the Third Council of Constantinople (681) defined that in Christ there are ``two natural wills\dots\ and two natural operations,'' not opposed to one another, but his human will following and freely subject to his divine and omnipotent will.\endnote{Third Council of Constantinople (681), Actio XVIII; quoted by Aquinas in \emph{ST} III, q.\ 18, a.\ 1, corpus and a.\ 6, sed contra.}

St.\ Thomas Aquinas analyzes the interior movements of Christ's human soul at Gethsemane with exacting care.\endnote{Aquinas, \emph{ST} III, q.\ 18, aa.\ 1, 4, 5, and 6.} Because Christ assumed a complete created human nature, he possessed a genuine human free will whose proper act is choice (\emph{ST} III, q.\ 18, a.\ 4). Within that human nature, Aquinas distinguishes three movements: the sensitive appetite (which naturally recoils from bodily pain), the rational human will considered as nature (\emph{voluntas ut natura}, which regards things in themselves and naturally turns away from death as an evil contrary to human nature), and the rational human will considered as reason (\emph{voluntas ut ratio}, which deliberates and chooses means in relation to an end).

Before the Passion, by a free dispensation of love, the Son of God allowed his passible flesh and all the powers of his human soul to experience what naturally belongs to them (\emph{ST} III, q.\ 18, a.\ 5). Therefore:

\begin{quote}
it is clear that the will of sensuality naturally shrinks from sensible pains and bodily hurt. In like manner, the will as nature turns from what is against nature and what is evil in itself, as death and the like; yet the will as reason may at time choose these things in relation to an end\dots\ Now it was the will of God that Christ should undergo pain, suffering, and death, not that these of themselves were willed by God, but for the sake of man's salvation. Hence it is plain that in His will of sensuality and in His rational will considered as nature, Christ could will what God did not; but in His will as reason He always willed the same as God, which appears from what He says (Matt.\ 26:39): ``Not as I will, but as Thou wilt.''\endnote{Aquinas, \emph{ST} III, q.\ 18, a.\ 5, corpus (English Dominican translation, Project Gutenberg eBook 19950).}
\end{quote}

There was no sin, no rebellion, and no true contrariety of wills in Christ (\emph{ST} III, q.\ 18, a.\ 6), because neither his sensitive horror nor his natural recoil from death ever rejected the end of human salvation or held back his deliberative choice to drink the chalice.

Gethsemane settles three questions for every creature before God. First, it refutes the morbid notion that pain or death is good in itself, or that holiness forbids horror in the face of anguish: Christ's sinless human nature recoiled so violently from torture and death that he sweat blood and asked, if possible, that the chalice pass. Second, it vindicates the dignity of a return offered when sensibility is torn apart rather than consoled: if the incarnate Son himself prayed ``not as I will but as thou wilt'' in darkness and agony, no believer need mistake dryness or dread for insincerity. Third, it reveals perfect creaturely obedience from within---not the erasure of human freedom, but a true created human will gathering its terrified, mortal humanity and handing it freely over to the Father in love.

In that prayer, in the face of sin and death, a human offering matches what God deserves---not because a mere creature manufactured equivalence out of dust, but because the Person offering that genuine created human will is the eternal Son himself. Every Eucharistic chalice the Church lifts, every imperfect act of Sunday worship, every reluctant step into the vineyard, and every hidden work of mercy toward the poor is accepted because it is joined to that human \emph{fiat} uttered in Gethsemane and consummated on the Cross.

\subsection{Broken Bonds and the Sacrifice of Praise}

In the light of Gethsemane, the closing verses of Psalm 115 reveal their full weight. Having asked ``What shall I render to the Lord, for all the things that he hath rendered to me?'' (Ps 115:12) and answered by taking the chalice of salvation (Ps 115:13), the psalmist arrives at his final confession:

\begin{quote}
O Lord, for I am thy servant: I am thy servant, and the son of thy handmaid. Thou hast broken my bonds: I will sacrifice to thee the sacrifice of praise, and I will call upon the name of the Lord. (Ps 115:16--17; Heb.\ 116:16--17)
\end{quote}

The psalmist calls himself God's servant in the very breath in which he rejoices that God has shattered his chains: \emph{Dirupisti vincula mea}---``Thou hast broken my bonds.'' St.\ Augustine, reaching this verse in his \emph{Exposition on Psalm 115}, sees why freedom and belonging do not cancel each other:

\begin{quote}
For I have not found any deserts of mine, since Thou hast broken my bonds asunder; I therefore owe Thee the sacrifice of praise; because, although I will boast that I am Thy servant, and the son of Thy handmaid, I will glory not in myself, but in Thee, my Lord, who hast broken asunder my bonds, that when I return from my desertion, I may again be bound unto Thee.\endnote{Augustine, \emph{Enarrationes in Psalmos}, Ps 115, \S 12 (NPNF I.8, numbered Ps CXVI).}
\end{quote}

God breaks the bonds of sin, death, and servile fear so that the creature may be bound again to him in freedom and praise---rather than chained by ingratitude to creatures that cannot sustain him.

\subsection{Why the Account Stays Open}

If existence were a loan extended by a stranger, closing the account would be the goal of honor: you would work until principal and interest were retired, obtain a receipt, and owe the lender nothing more. Before God, that dream is both metaphysically impossible and spiritually ruinous. Metaphysically impossible, because every breath by which we pray, every coin placed in the hands of the poor, and every impulse of grace by which we repent and persevere is itself freshly given by God; even our merits are God's own gifts crowned by his promise, and gratitude grows rather than shrinks the more we give thanks. Spiritually ruinous, because closing the account with God would mean terminating our dependence on the Source of life---another name for death.

An open account before a tyrant is terror. An open account before the Father who gave his Son to drink our chalice is home. Justice without equivalence brings the creature honestly to his knees, teaches his body and lips the common speech of the Church, sends his hands outward to the poor, and guards him through repentance and moderation against both negligence and pride. Yet justice, even at its highest peak in the virtue of religion, is never the last word between God and the soul. Religion orders our acts toward God as our Creator and Sovereign; it stands waiting to be taken up, animated, and surpassed by the theological charity in which God no longer deals with us merely as servants rendering what is due, but calls us friends. Where every calculation of debt finally yields to love that asks no reason beyond the Beloved himself, the open account ceases to look like a ledger at all, and becomes communion without end.
